\section{Vocabulary List: Czech A1}

% 1. Nouns
\subsection{Nouns (Substantiva)}
Nouns in Czech are categorized into three genders: Masculine, Feminine, and Neuter. Below are common nouns grouped by gender.

\begin{center}
\small
\begin{tabular}{ll | ll | ll}
\hline
\multicolumn{2}{c|}{\textbf{Masculine (Ten)}} & \multicolumn{2}{c|}{\textbf{Feminine (Ta)}} & \multicolumn{2}{c}{\textbf{Neuter (To)}} \\ \hline
\textbf{Czech} & \textbf{English} & \textbf{Czech} & \textbf{English} & \textbf{Czech} & \textbf{English} \\ \hline
Dům & House & Žena & Woman & Auto & Car \\
Svetr & Sweater & Káva & Coffee & Divadlo & Theater \\
Čaj & Tea & Škola & School & Německo & Germany \\
Grep & Grapefruit & Tramvaj & Tram (Exc.) & Město & City \\
Guláš & Goulash & Restaurace & Restaurant & Pivo & Beer \\
Hrad & Castle & Praha & Prague & Nádraží & Station \\
Oběd & Lunch & Banka & Bank & Jídlo & Food/Meal \\
Film & Movie & Ulice & Street & Letiště & Airport \\
Hotel & Hotel & Rodina & Family & Kino & Cinema \\
Mobil & Mobile phone & Práce & Work/Job & Moře & Sea \\ \hline
\end{tabular}
\end{center}

% 2. Verbs
\subsection{Verbs (Slovesa)}
These are the core verbs from the "MÁME NOVÝ TEPLÝ SVETR" exercises, organized by conjugation type.

\begin{center}
\small
\begin{tabular}{lll | lll}
\hline
\textbf{Infinitive} & \textbf{English} & \textbf{Type} & \textbf{Infinitive} & \textbf{English} & \textbf{Type} \\ \hline
Být & To be & Irregular & Vařit & To cook & -Í- Type \\
Jíst & To eat & Irregular & Učit se & To learn/study & -Í- Type \\
Dělat & To do/make & -Á- Type & Rozumět & To understand & -Í- Type \\
Mít & To have & -Á- Type & Chodit & To go (foot) & -Í- Type \\
Obědvat & To lunch & -Á- Type & Jezdit & To go (vehicle) & -Í- Type \\
Běhat & To run & -Á- Type & Pracovat & To work & -UJ- Type \\
Dívat se & To watch & -Á- Type & Studovat & To study & -UJ- Type \\
Poslouchat & To listen & -Á- Type & Nakupovat & To shop & -UJ- Type \\
Mluvit & To speak & -Í- Type & Pít & To drink & -U- Type \\
Číst & To read & -Í- Type* & Psát & To write & -U- Type \\ \hline
\end{tabular}
\end{center}

% 3. Adjectives
\subsection{Adjectives (Adjektiva)}
Hard adjectives change their endings based on gender: \textbf{-ý} (M), \textbf{-á} (F), \textbf{-é} (N). Soft adjectives (e.g., \textit{moderní}) remain the same for all genders.

\begin{center}
\small
\begin{tabular}{lll | lll}
\hline
\textbf{Czech (M)} & \textbf{English} & \textbf{Opposite} & \textbf{Czech (M)} & \textbf{English} & \textbf{Opposite} \\ \hline
Dobrý & Good & Špatný (Bad) & Hezký & Nice/Pretty & Ošklivý (Ugly) \\
Velký & Big & Malý (Small) & Rychlý & Fast & Pomalý (Slow) \\
Nový & New & Starý (Old) & Levný & Cheap & Drahý (Expensive) \\
Teplý & Warm & Studený (Cold) & Šťastný & Happy & Smutný (Sad) \\
Zajímavý & Interesting & Nudný (Boring) & Vysoký & Tall/High & Malý/Nízký \\
Moderní & Modern & (Soft Adj.) & Inteligentní & Intelligent & Hloupý (Stupid) \\
Čerstvý & Fresh & Starý (Stale) & Hubený & Thin & Tlustý (Fat) \\ \hline
\end{tabular}
\end{center}
\small \textbf{Grammar Note:} 
Remember gender agreement: \textit{Ten dům je velk\textbf{ý}}, \textit{Ta káva je dobr\textbf{á}}, \textit{To auto je nov\textbf{é}}.

% 4. Colors (Barvy)
\subsection{Colors (Barvy)}
Colors are "hard adjectives."\\
They must match the gender of the noun:
\textbf{Masculine: -ý} | \textbf{Feminine: -á} | \textbf{Neuter: -é}.

\begin{center}
\begin{tabular}{ll|ll|ll}
\hline
\textbf{Czech} & \textbf{English} & \textbf{Czech} & \textbf{English} & \textbf{Czech} & \textbf{English} \\ \hline
Bílý    & White  & Černý    & Black  & Červený  & Red    \\
Modrý   & Blue   & Zelený   & Green  & Žlutý    & Yellow \\
Hnědý   & Brown  & Šedý     & Grey   & Oranžový & Orange \\
Růžový  & Pink   & Fialový  & Purple & Světlý   & Light  \\ \hline
\end{tabular}
\end{center}